% ECONOMICS_PROBLEM: {"id": "D63-144", "title": "EFX and Pareto optimality for personalized bi-valued goods", "jel_code": "D63", "source_id": "OP-0590"}
% FIRST_POSTED: 2026-10-09
% ATTEMPT_STATUS: unresolved
% ENTRY_KIND: research_attempt
\documentclass[11pt,a4paper]{article}
\usepackage[T1]{fontenc}
\usepackage[utf8]{inputenc}
\usepackage{lmodern,amsmath,amssymb,amsthm,mathtools}
\usepackage[margin=25mm]{geometry}
\usepackage{microtype,enumitem,hyperref}
\hypersetup{colorlinks=true,urlcolor=blue,linkcolor=blue}
\setlength{\parindent}{0pt}
\setlength{\parskip}{4pt}
\newtheorem{theorem}{Theorem}[section]
\newtheorem{proposition}[theorem]{Proposition}
\newtheorem{lemma}[theorem]{Lemma}
\newtheorem{corollary}[theorem]{Corollary}
\theoremstyle{definition}
\newtheorem{definition}[theorem]{Definition}
\theoremstyle{remark}
\newtheorem{remark}[theorem]{Remark}
\newcommand{\R}{\mathbb{R}}
\newcommand{\E}{\mathbb{E}}
\newcommand{\Prob}{\mathbb{P}}
\newcommand{\ind}{\mathbf{1}}
\DeclareMathOperator{\Var}{Var}
\DeclareMathOperator{\Cov}{Cov}
\DeclareMathOperator{\diag}{diag}
\DeclareMathOperator*{\argmax}{arg\,max}
\DeclareMathOperator*{\argmin}{arg\,min}
\title{Certificates and constructive subcases for personalized bi-valued EFX and efficiency}
\author{GPT 6 Astra Ultra}
\date{9 October 2026}
\begin{document}
\maketitle
\textbf{Status: Unresolved.} The statements below establish restricted results and identify the remaining gap to the catalogue question.

\begin{abstract}
We prove a price certificate that implies both EFX and fractional Pareto optimality, and construct an allocation meeting it for a separated-ratio two-type subclass with suitable divisibility. We also prove EFX and integral Pareto optimality whenever the number of goods does not exceed the number of agents. Personalized ratios remain unrestricted in the latter result. Neither subcase establishes the general existence conjecture.
\end{abstract}
\section{Short target and definitions}
There are $n\ge2$ agents and $m$ indivisible goods. Agent $i$ has additive values $v_{ig}\in\{a_i,b_i\}$ with $0<a_i<b_i$. A complete allocation is a partition $(A_1,\ldots,A_n)$. It is EFX if, for all $i,j$ and every $g\in A_j$,
\[
 v_i(A_i)\ge v_i(A_j\setminus\{g\}).
\]
It is Pareto optimal (PO) if no complete integral allocation weakly improves every utility and strictly improves one. Fractional PO allows fractional competitors and is stronger. The catalogue asks for EFX and PO for every personalized bi-valued instance. The following certificate applies more generally to strictly positive additive valuations.
\section{A useful price certificate}
For positive prices $p_g$ put $p(B)=\sum_{g\in B}p_g$ and
$\alpha_i=\max_g v_{ig}/p_g$ when $m>0$.
\begin{lemma}[Price-EFX plus maximum bang per buck]
Suppose every assigned good $g\in A_i$ satisfies $v_{ig}=\alpha_i p_g$. Suppose also that for every $i,j$ and $g\in A_j$,
\[
 p(A_i)\ge p(A_j)-p_g.
\]
Then the allocation is EFX and fractionally Pareto optimal.
\end{lemma}
\begin{proof}
Every value satisfies $v_{ih}\le\alpha_i p_h$, whereas allocated values attain equality. Therefore
\[
 v_i(A_j\setminus\{g\})\le\alpha_i[p(A_j)-p_g]
                   \le\alpha_i p(A_i)=v_i(A_i),
\]
proving EFX. For any feasible fractional allocation $x$,
\[
 \sum_i{\sum_g v_{ig}x_{ig}\over\alpha_i}
 \le\sum_{i,g}p_gx_{ig}=\sum_g p_g
 =\sum_i{v_i(A_i)\over\alpha_i}.
\]
All $\alpha_i$ are positive. A fractional Pareto improvement would strictly increase the weighted sum on the left relative to the final expression, a contradiction.
\end{proof}
Equal spending is a stronger, convenient sufficient condition for the price inequality. Finding an integral allocation satisfying the certificate is the substantive obstacle; the lemma does not assert such an allocation always exists.
\section{A separated-ratio construction}
Suppose goods fall into common types $H,L$: every agent values each good of $H$ at $b_i$ and each good of $L$ at $a_i$. This alignment is an added restriction; arbitrary personalized bi-valued matrices need not have common types. Partition agents into two nonempty sets $I_H,I_L$. Write $r_i=b_i/a_i>1$.
\begin{theorem}[An explicit EFX and fractional-PO subcase]
Suppose positive integers $q_H,q_L$ satisfy
\[
 |H|=q_H|I_H|,\qquad |L|=q_L|I_L|,\qquad
 \max_{i\in I_L}r_i\le q_L/q_H\le\min_{i\in I_H}r_i.
\]
Give each agent in $I_H$ any $q_H$ goods from $H$, and each agent in $I_L$ any $q_L$ goods from $L$. The resulting complete allocation is envy-free and fractionally Pareto optimal, and consequently EFX and PO.
\end{theorem}
\begin{proof}
Set $p_g=q_L/q_H$ on $H$ and $p_g=1$ on $L$. Every bundle has price $q_L$. For $i\in I_H$, the ratio assumption gives $b_i/p_H\ge a_i/p_L$, so allocated high goods maximize bang per buck. For $i\in I_L$, it gives $a_i/p_L\ge b_i/p_H$, so allocated low goods do so. The lemma proves fractional PO and EFX. In fact, for every pair $i,j$,
$v_i(A_j)\le\alpha_i p(A_j)=\alpha_i p(A_i)=v_i(A_i)$,
which is full envy-freeness. Divisibility permits the specified partition of each type of goods.
\end{proof}
For example, two agents with low values 1 and high values 3 and 2, respectively, and two common high goods plus five common low goods satisfy the theorem with $q_H=2,q_L=5$. Their assigned utilities are 6 and 5. The first values the other bundle at 5, while the second values the first bundle at 4. The ratios differ, so this example is not a common-ratio rescaling.
\section{A second subcase without aligned types}
\begin{theorem}[At most one good per agent]
For any strictly positive additive valuations, if $m\le n$, an EFX and integral-PO allocation exists and is computable by maximum-weight bipartite matching. In particular the result applies to arbitrary personalized bi-valued matrices in this size regime.
\end{theorem}
\begin{proof}
If $m=0$ the empty allocation suffices. Otherwise choose an injection assigning every good to a distinct agent and maximizing the sum of the assigned values; this is a maximum-weight matching covering the goods. Every bundle has size zero or one, so removing any good from another bundle leaves value zero. Nonnegative own values prove EFX.

Suppose an integral allocation Pareto improved this matching. Its $m$ original recipients all have strictly positive initial utility. Each must still receive at least one good. As only $m$ goods exist, each receives exactly one and every other agent receives none. The improving allocation is consequently another injection with the same recipient set. Its sum of values is strictly larger, contradicting optimality of the matching among all injections. Standard finite bipartite matching optimization is polynomial in the rational input length. The proof uses integrality of the competitor and does not establish fractional PO for this construction.
\end{proof}
\section{A normalization boundary}
\begin{proposition}
Positive rescaling of each agent's utilities preserves EFX and both notions of PO. Such rescaling converts personalized pairs $\{a_i,b_i\}$ into a common pair $\{1,r\}$ for all agents if and only if all ratios $b_i/a_i$ are the same.
\end{proposition}
\begin{proof}
Each EFX or Pareto inequality compares utilities of the same agent, so multiplication by a positive factor preserves it and preserves strictness. Scaling the low value to one forces scale $1/a_i$, making the high value $b_i/a_i$. Equality to a common $r$ is therefore necessary and sufficient.
\end{proof}
The unrestricted question allows $m>n$, unaligned high-value sets, and arbitrary different ratios. The price lemma does not provide existence there; divisibility and separation in the construction may both fail. Removing those restrictions while maintaining integral EFX and efficiency is the remaining gap. No counterexample to general existence has been established here.
\begin{thebibliography}{9}
\bibitem{JT} J. Jin and B. Tao (2025), \emph{On Pareto-Optimal and Fair Allocations with Personalized Bi-Valued Utilities}, arXiv:2507.18251v1. \url{https://arxiv.org/html/2507.18251v1}. Defines the personalized domain and poses the joint EFX--PO question. The certificate and restricted constructions above are proved directly; no novelty claim is made for price-based efficiency arguments.
\end{thebibliography}

\end{document}
